% Lösungen Zone (zusammenbau v0.1)
\einheitenkopf[][Kennst du schon]{Das kennst du schon}

Zuordnung: 1 → Zähler und Nenner der Formel · 2 → der zweite Weg und die Mindestens-Fälle · 3 → die Modellwahl · 4 → die Bruchkette · 5 → der zweite Weg und die Mindestens-Fälle · 6 → der zweite Weg und die Mindestens-Fälle

\erg{1}{a) $(5 \text{ über } 2) = 10$ \quad b) $(10 \text{ über } 3) = 120$}
\erg{2}{a) $\frac{3}{5} \cdot \frac{2}{4} = \frac{3}{10}$ \quad b) $\frac{4}{7} \cdot \frac{3}{6} \cdot \frac{2}{5} = \frac{24}{210} = \frac{4}{35}$}
\erg{3}{a) ja – zehn unabhängige Würfe mit immer gleicher Wahrscheinlichkeit $\frac{1}{6}$ für eine Sechs. \quad b) ja – bei $50\,000$ Einwohnern ändert das Herausgreifen von $20$ Personen den Anteil kaum.}
\erg{4}{a) $\frac{6}{12} = \frac{1}{2}$ \quad b) $\frac{60}{336} = \frac{5}{28}$}
\erg{5}{a) Fehler: Tim rechnet, als würde die Kugel zurückgelegt. Nach dem ersten Zug sind nur noch $2$ schwarze von $7$ Kugeln da. Richtig: $\frac{3}{8} \cdot \frac{2}{7} = \frac{6}{56} = \frac{3}{28}$}
\erg{6}{a) $\frac{2}{8} \cdot \frac{1}{7} = \frac{2}{56} = \frac{1}{28}$}
